
\needspace{4\baselineskip}
\color{uniblue}{\subsection{Контроль наличия напряжения на шинах (КННш)}}
\color{black}

\begin{enumerate}[label=\arabic{section}.\arabic{subsection}.\arabic*, labelsep=4pt, leftmargin=0pt, itemindent=57pt]

\item
Функция контролирует наличие напряжения на секции шин с дополнительным контролем симметрии по $U_2$.
% УСЛОВНОЕ ВЫРАЖЕНИЕ
\ifboolexpr{bool {isVV}}
{
    Функция используется для АПВ и АВР, а также для оперативного включения.
}
{
    Функция используется для АПВ, а также для оперативного включения.
}

\item
Пусковые условия функции формируются без выдержки времени при превышении всех линейных напряжений на шинах выше уставки и одновременном выполнении следующих условий, если уставка SGF2 выставлена в положение <<Измеренный>>:
\begin{itemize}
    \item функция введена в работу;
    \item минимальное из подведенных линейных напряжений, измеренных с ТН шин, выше уставки;
    \item отсутствует срабатывание ИО напряжения обратной последовательности.
\end{itemize}
\item
Предусмотрена блокировка функции при срабатывании ИО обратной последовательности напряжения выше уставки.
\item
Если уставка SGF2 выставлена в положение <<Внешний>>, работа функции осуществляется при наличии активного внешнего дискретного сигнала <<Наличие Uш>> без учета срабатывания измерительных органов.

\item
Ввод функции в работу на этапе параметрирования устройства осуществляется программным переключателем SGF1 <<Ввод\_функции>>. Информация о введенной функции отображается выходным сигналом <<КННш: Ввод>>.

\item
Функция <<КННш>> может быть выведена из работы оперативно путем активации сигнала <<ОВ КННш>>, что характеризуется активным состоянием сигнала <<КННш: Оперативный вывод>> на выходе функции. 
\item
Алгоритм работы функции <<КННш>> выполнен в соответствии с рисунком \ref{knnsh:img1}. В таблице \ref{knnsh:tbl1} приведены параметры, необходимые для настройки функции <<КННш>>.

\vspace{3mm}
\begin{figure}[H]
\centering
\includegraphics[width=0.9\textwidth,height=0.9\textheight,keepaspectratio]{\fbpath/122. КННш/_latex/img/knnsh.pdf}
\captionof{figure}{Функциональная схема <<КННш>>}\label{knnsh:img1}
\end{figure}

\small
\begin{longtable}{|>{\centering\arraybackslash}m{5.3cm}|>{\centering\arraybackslash}m{3.3cm}|>{\centering\arraybackslash}m{4.2cm}|>{\centering\arraybackslash}m{1.8cm}|>{\centering\arraybackslash}m{1cm}|}
\caption{Параметры для настройки функции <<КННш>>\hfill\vspace{-0.5\baselineskip}}\label{knnsh:tbl1}\\ 
\hline
\rowcolor{gray!30}
Параметр (Параметр на ИЧМ) & Условное обозначение на схеме & Значение/ Диапазон & Единица измерения & Шаг \\ 
\hline
\endfirsthead
\caption*{\hspace{3pt}\emph{Продолжение таблицы \ref{knnsh:tbl1}\hfill\vspace{-0.5\baselineskip}}} \\ % сделано по ГОСТ 2.105 п.6.8.7
\hline
\rowcolor{gray!30}
Параметр (Параметр на ИЧМ) & Условное обозначение на схеме & Значение/ Диапазон & Единица измерения & Шаг \\ 
%\hline
\endhead
%\hline
\endfoot
%\hline
\endlastfoot
%==+t1*PTOV1|VOLPRES1
\centering Ввод функции в работу (Ввод\_функции) & \centering SGF1 & \centering 0 = Не предусмотрено\\1 = Предусмотрено & \centering -- & \centering \arraybackslash -- \\
\hline
\centering Напряжение срабатывания (Uср) & \centering U> & \centering 50,0 ... 150,0 & \centering В & \centering \arraybackslash 0,1 \\
\hline
\centering Напряжение срабатывания ОП (U2ср) & \centering U2> & \centering 2,0 ... 50,0 & \centering В & \centering \arraybackslash 0,1 \\
\hline
\centering Режим пуска (Реж\_пуска) & \centering SGF2 & \centering 0 = Измеренный\\1 = Внешний & \centering -- & \centering \arraybackslash -- \\
\hline
%===t1
\end{longtable}
\normalsize

\end{enumerate}
